\section{Retain an incomplete block between updates}
The streaming context stores the information that must survive a return to the caller.
Its eight state words describe all complete message blocks processed so far. Its byte
count includes every original message byte supplied so far. Its buffer contains the
remaining bytes, which are fewer than 64 after an update returns.

\begin{table}[htbp]
\centering\caption{The project's context layout. Offsets and sizes are decimal bytes.}
\tagpdfsetup{table/header-rows={1}}
\begin{tabular}{rrl}
\toprule Offset & Size & Meaning\\\midrule
0 & 32 & Eight numeric state words\\
32 & 8 & Original message byte count\\
40 & 4 & Number of occupied buffer bytes, zero through 63\\
44 & 4 & Reserved context storage\\
48 & 64 & Incomplete block buffer\\
112 & 256 & Schedule workspace\\
\bottomrule
\end{tabular}
\end{table}

Initialization copies the eight initial state words, sets the byte count to zero and
sets the used count to zero. Old buffer and workspace contents do not need to be
cleared: their contents have no meaning until the program writes the bytes or words it
will use. This saves initialization work, but it also makes the used count essential.
A zero-filled allocation must not conceal a missing buffer write.

On update, add the incoming length to the 64-bit byte count. Copy source bytes to the
buffer beginning at the used count. When it reaches 64, compress that buffer and set
the used count back to zero. Continue copying until the selected source interval is
exhausted. Store the remaining used count before returning. The implementation here
copies all input through the buffer; processing aligned full blocks directly from the
source is a possible later optimization.

Consider a 65-byte message supplied as calls of 55 and 10 bytes. After the first call,
no full block has been processed, the total is 55 and the used count is 55. The next
nine bytes complete the block. Compress it, clear the used count, then copy the last
byte to buffer position zero. On return, the total is 65 and the used count is one.
A single 65-byte update must leave the same meaningful context state.

This observation provides a test before finalization exists. Compare the eight state
words after the complete block, the total, the used count and the remaining buffer
prefix. Bytes beyond that prefix are unspecified workspace. Requiring them to match
would reject correct implementations for an irrelevant storage choice.
